\section{Introduction}
Trentino has a dense network of official hiking trails, numbered, signposted and maintained by the \emph{Società degli Alpinisti Tridentini} (SAT)\footnote{\url{https://www.sat.tn.it}}, the Trentino section of the Italian Alpine Club. Hiking is one of the main outdoor activities for both residents and visitors, from short walks to a mountain hut to demanding high-altitude routes.

\noindent However, the information needed to plan a hike is scattered across many sources. Trail descriptions (route, length, duration, difficulty) are kept in the SAT trail register and in OpenStreetMap (OSM)\footnote{\url{https://www.openstreetmap.org}}; mountain huts, restaurants and car parks in other datasets; public transport in the timetables of Trentino Trasporti\footnote{\url{https://www.trentinotrasporti.it}}; and current conditions, such as recent rain, snow cover and the forecast, only in the data of the provincial weather service, Meteotrentino\footnote{\url{https://www.meteotrentino.it}}. In addition, the official difficulty scale of the \emph{Club Alpino Italiano} (CAI), with its levels T (tourist), E (hiking), EE (expert hikers) and EEA (expert hikers with equipment), is meaningful to experienced hikers but not to beginners, who are the ones most in need of guidance.

\noindent A Knowledge Graph (KG) \cite{p4,p6} can integrate this heterogeneous information into one connected structure that can be queried as a whole. By linking trails with their starting points, nearby car parks and public transport stops, huts and restaurants, lifts, weather observations and forecasts, and rules about equipment and clothing, the KG can answer the practical questions a hiker asks before setting out. This project builds such a KG for Trentino following the iTelos methodology \cite{p1,p2,p3,p5}, in which the purpose of the KG, expressed through scenarios, personas and competency questions, drives every subsequent design decision.

\noindent The report is structured as follows:
\begin{itemize}
    \item Section 2: the project design, i.e.\ the domain of interest and the informal purpose.
    \item Sections 3, 4, 5, 6: the iTelos phases (purpose, language, knowledge and entity definition) and their activities, divided by knowledge and data layer, with the qualitative evaluation of the resources produced.
    \item Sections 7 and 8: the quantitative evaluation and the metadata of the resources handled and generated by the project.
    \item Section 9: conclusions and open issues.
\end{itemize}
